%
% reversi_020_introduzione.tex
%
% This file is part of the reversi program
% http://github.com/rcrr/reversi-v2
%
% Author Roberto Corradini mailto:rob_corradini@yahoo.it
% Copyright 2026 Roberto Corradini. All rights reserved.
%
% License
%
% This program is free software; you can redistribute it and/or modify it
% under the terms of the GNU General Public License as published by the
% Free Software Foundation; either version 3, or (at your option) any
% later version.
%
% This program is distributed in the hope that it will be useful,
% but WITHOUT ANY WARRANTY; without even the implied warranty of
% MERCHANTABILITY or FITNESS FOR A PARTICULAR PURPOSE. See the
% GNU General Public License for more details.
%
% You should have received a copy of the GNU General Public License
% along with this program; if not, write to the Free Software
% Foundation, Inc., 59 Temple Place - Suite 330, Boston, MA  02111-1307, USA
% or visit the site <http://www.gnu.org/licenses/>.
%
\chapter{Introduzione}

Reversi è un ecosistema di programmi e strumenti informatici per studiare il gioco da tavolo Othello, anche chiamato Reversi.
L'intento è studiare, capire, sviluppare algoritmi e modelli che possano progredire nella realizzazione di un agente
sintetico ed autonomo che sia in grado di giocare una partita in maniera autorevole. E in prospettiva di tendere alla
risoluzione matematica del gioco, o quantomeno ad una sua approssimazione probabilistica.
I programmi sono scritti principalmente in ANSI C, Python, SQL, con altri linguaggi come Common Lisp, Java, Clojure, R usati
per sperimentazioni, ma non rilevanti nella sostanza.
Il corpo principale del programma è sviluppato in C, e comprende l'endgame solver, il generatore di partite random, il
solutore integrato con il database REGAB, e le funzioni di risoluzione del modello di ML basato sulla Regressione Logistica.
Il Database SQL REGAB, sviluppato in PostgreSQL, ed in linguaggio procedurale PL/pgSQL, contiene le posizioni random
che poi vengono risolte e taggate con il valore vero del game.
Il codice in C è stato poi portato in Python, per poter operare più agilmente con l'ecosistema attivo del Machine-Learning, AI,
calcolo scientifico, acceleratori GPU, e tutto quanto serve per la ricerca di frontiera.
Il codice in C rimane più performante per la risoluzione delle posizioni, il codice in Python è più consono allo sviluppo
dei Modelli per la definizione di evaluation function più evolute.
